%%%%%%%%%%%%%%%%%%%%%%%%%%%%%%%%%%%%%%%%%%%%%%%%%%%%%%%%%%%%%%%%%%%%%%%%%%%%%%%%
% Preámbulo                                                                    %
%%%%%%%%%%%%%%%%%%%%%%%%%%%%%%%%%%%%%%%%%%%%%%%%%%%%%%%%%%%%%%%%%%%%%%%%%%%%%%%%

\documentclass[11pt,a4paper,titlepage,oneside]{report}

%%% RELACIÓN DE VARIABLES A PERSONALIZAR %%%
%\def\lingua{gal}
\def\lingua{esp} % descomenta esta liña se redactarás a memoria en español
%\def\lingua{eng} % descomenta esta liña se redactarás a memoria en inglés
\def\nome{MIRYAM RAQUEL SÁNCHEZ TORRE, YAGO FERNÁNDEZ REGO, ROSANA LLILLART GUIBERT} % substitúe aquí o teu nome
\def\nomeA{MIRYAM RAQUEL SÁNCHEZ TORRE} 
\def\nomeB{YAGO FERNÁNDEZ REGO} 
\def\nomeC{ROSANA LLILLART GUIBERT}   
                          % substitúe aquí o teu nome
\def\nomedirectorA{JAVIER PARAPAR}              % substitúe aquí o nome de quen dirixe
\def\nomedirectorB{ANXO PÉREZ}             % duplica esta liña máis veces se o precisas, cambiando
                                                     % a letra final (A, B, C, D...): úsanse na portada.tex
\def\titulo{MYR:Buscador de viajes por España} % substitúe aquí o título do teu TFG
%\def\titulacion{gced}                               % descomenta esta liña e comenta a seguinte se es estudante do GCED
\def\titulacion{MUEI}
%\def\titulacion{gia}                                % comenta a liña anterior e descomenta esta se es estudante do GIA
%\def\mencion{COMPUTACIÓN}                           % descomenta a mención que che corresponda se es estudante do GEI
\def\mencion{MASTER DE INGENIERÍA INFORMÁTICA} % descomenta a mención que che corresponda se es estudante do GEI
%\def\mencion{ENXEÑARÍA DE COMPUTADORES}
%\def\mencion{SISTEMAS DE INFORMACIÓN}
%\def\mencion{TECNOLOXÍAS DA INFORMACIÓN}

%\def\renomearcadros{si} % descomenta esta liña se redactas a memoria en español e prefires que
                         % os "cuadros" e o "índice de cuadros" se renomeen
                         % a "tablas" e "índice de tablas" respectivamente

\usepackage{estilo_tfg}

% Lista de paquetes potencialmente interesantes (uso baixo demanda)

% \usepackage{alltt}       % proporciona o entorno alltt, semellante a verbatim pero que respecta comandos
% \usepackage{enumitem}    % permite personalizar os entornos de lista
% \usepackage{eurofont}    % proporciona o comando \euro
% \usepackage{float}       % permite máis opcións para controlar obxectos flotantes (táboas, figuras)
% \usepackage{hhline}      % permite personalizar as liñas horizontais en arrays e táboas
  \usepackage{longtable}   % permite construir táboas que ocupan máis dunha páxina
% \usepackage{lscape}      % permite colocar partes do documento en orientación apaisada
% \usepackage{moreverb}    % permite personalizar o entorno verbatim
  \usepackage{multirow}    % permite crear celdas que ocupan varias filas da mesma táboa
% \usepackage{pdfpages}    % permite insertar ficheiros en PDF no documento
% \usepackage{rotating}    % permite diferentes tipos de rotacións para figuras e táboas
% \usepackage{subcaption}  % permite a inclusión de varias subfiguras nunha figura
% \usepackage{tabu}        % permite táboas flexibles
% \usepackage{tabularx}    % permite táboas con columnas de anchura determinada

%%%%%%%%%%%%%%%%%%%%%%%%%%%%%%%%%%%%%%%%%%%%%%%%%%%%%%%%%%%%%%%%%%%%%%%%%%%%%%%%
% Corpo                                                                        %
%%%%%%%%%%%%%%%%%%%%%%%%%%%%%%%%%%%%%%%%%%%%%%%%%%%%%%%%%%%%%%%%%%%%%%%%%%%%%%%%

\begin{document}

 %%%%%%%%%%%%%%%%%%%%%%%%%%%%%%%%%%%%%%%%
 % Preliminares do documento            %
 %%%%%%%%%%%%%%%%%%%%%%%%%%%%%%%%%%%%%%%%

 \begin{titlepage}
  
  \hspace*{128pt}
  \textcolor{udcpink}{{\fontencoding{T1}\fontfamily{phv}\selectfont Facultad de Informática}}\\[-32pt]

  \begin{center}
    \includegraphics[scale=0.3]{imaxes/udc}\\[25pt]

    {\large TRABAJO FIN DE CURSO \\
            \nometitulacion \\
            \nomemencion } \\[10pt]

    \carimbo \\[25pt]

    \begin{huge}
      \begin{spacing}{1.3}
        \bfseries \titulo
      \end{spacing}
    \end{huge}
  \end{center}
  
  \vfill
  
  \begin{flushright}
    {\large
    \begin{tabular}{ll}
      {\bf Estudante:} & \nomeA\\
                       & \nomeB\\ 
                       & \nomeC\\

      
      {\bf Dirección:} & \nomedirectorA \\
                       & \nomedirectorB \\ % duplica esta liña máis veces se o precisas, cambiando
                                           % a letra final (A, B, C, D...); define eses nomes no memoria_tfg.tex
    \end{tabular}}
  \end{flushright}
  \rightline{A Coruña, \datasimple.}
\end{titlepage}

%\dedicatoria{Dedicatoria} % escribe neste comando o teu texto de dedicatoria
%\paxinaenbranco
%\begin{agradecementos}
% \blindtext                % substitúe este comando polo teu texto de agradecementos
 %\end{agradecementos}
 %\include{portada/resumo}

 \pagenumbering{roman}
 \setcounter{page}{1}
 %\bstctlcite{IEEEexample:BSTcontrol}

 \tableofcontents
 %\listoffigures
 %\listoftables
 \clearpage
 
 \pagenumbering{arabic}
 \setcounter{page}{1}

 %%%%%%%%%%%%%%%%%%%%%%%%%%%%%%%%%%%%%%%%
 % Capítulos                            %
 %%%%%%%%%%%%%%%%%%%%%%%%%%%%%%%%%%%%%%%%

%\include{contido/introducion}
%\chapter{Contido demostrativo}
\label{chap:demo}

\lettrine{E}{ntre} a introdución e as conclusións, o documento conterá
tantos capítulos como sexa preciso, sempre con coidado de non rebasar
o límite de 80 páxinas fixado polo regulamento de TFGs.

Empregaremos éste de xeito demostrativo, para ilustrar o uso de
elementos habituais que poidan ser de utilidade\footnote{Por exemplo,
  isto é unha nota a pé de páxina.}.

\section{Inclusión de imaxes}

Se precisamos imaxes no noso documento, incluirémolas do xeito que se
indica na figura~\ref{fig:exemplo} (páxina~\pageref{fig:exemplo}). Se
o facemos así, \LaTeX ubicará cada imaxe no mellor lugar posible,
lugar que pode variar a medida que o documento vaia crecendo coa
inclusión de máis texto e outros elementos (máis imaxes, táboas,
etc.).

\begin{figure}[hp!]
  \centering
  \includegraphics[width=0.75\textwidth]{imaxes/udc.png}
  \caption{Pé de imaxe descritivo}
  \label{fig:exemplo}
\end{figure}

Recoméndase almacenar os ficheiros gráficos no directorio
\texttt{imaxes}.

\subsection{Inclusión de varias sub-imaxes}

Se precisamos inserir imaxes relacionadas, pode ser apropiado
incluílas como sub-figuras, do xeito que se pode apreciar na
figura~\ref{fig:exemplo-subfiguras} (páxina~\pageref{fig:exemplo-subfiguras})
coas imaxes~\ref{fig:subfigura-rotada} e~\ref{fig:subfigura-deformada}.
Como se pode ver nos exemplos desta sección, sempre é recomendable
referirse ás imaxes (ou táboas e outros elementos \emph{flotantes},
que se demostrarán nas seccións seguintes deste capítulo demostrativo)
pola súa referencia, xa que dese xeito non dependemos de onde
queden ubicados os elementos en cuestión.

\begin{figure}[hp!]
  \centering
  \begin{subfigure}[c]{0.3\textwidth}
    \includegraphics[angle=45,width=\textwidth]{imaxes/udc.png}
    \caption{Pé de subimaxe rotada}
    \label{fig:subfigura-rotada}
  \end{subfigure}
  \hspace{0.1\textwidth}
  \begin{subfigure}[c]{0.3\textwidth}
    \includegraphics[width=\textwidth,height=3cm]{imaxes/udc.png}
    \caption{Pé de subimaxe deformada}
    \label{fig:subfigura-deformada}
  \end{subfigure}
  \caption{Pé de imaxe xeral}
  \label{fig:exemplo-subfiguras}
\end{figure}

\section{Inclusión de táboas}

Se precisamos táboas no noso documento, incluirémolas do xeito que se
indica na táboa~\ref{tab:exemplo} (páxina~\pageref{tab:exemplo}). Se
o facemos así, \LaTeX ubicará cada táboa no mellor lugar posible,
lugar que pode variar a medida que o documento vaia crecendo coa
inclusión de máis texto e outros elementos (máis imaxes, táboas,
etc.).

\begin{table}[hp!]
  \centering
  \rowcolors{2}{white}{udcgray!25}
  \begin{tabular}{c|c}
  \rowcolor{udcpink!25}
  \textbf{Título de columna} & \textbf{Outro título de columna} \\\hline
  \textit{Título de fila} & Contido da cela \\
  \textit{Título de fila} & Contido da cela \\
  \textit{Título de fila} & Contido da cela \\
  \textit{Título de fila} & Contido da cela \\
  \textit{Título de fila} & Contido da cela \\
  \textit{Título de fila} & Contido da cela \\
  \end{tabular}
  \caption{Pé de táboa descritivo}
  \label{tab:exemplo}
\end{table}

\subsection{Inclusión de táboas longas}

Para táboas longas que ocupan varias páxinas, como é o caso da \ref{tab:longa}
(páxina~\pageref{tab:longa}), recoméndase o uso do paquete \texttt{lontable},
incluído xa entre os paquetes recomendados no ficheiro raíz do proxecto
(\verb+memoria_tfg.tex+).

\input{contido/taboa-longa.tex}

\subsection{Inclusión de táboas con celas que ocupan varias columnas ou filas}

En ocasións pode resultar de interese incluír nunha táboa unha cela que se estenda
a través de varias columnas, como ocorre na táboa~\ref{tab:exemplocolumnas}
(páxina~\pageref{tab:exemplocolumnas}).

\begin{table}[hp!]
  \centering
  \rowcolors{2}{white}{udcgray!25}
  \begin{tabular}{c|c|c}
  \rowcolor{udcpink!25}
  \multicolumn{3}{c}{\textbf{Cela en varias columnas}} \\\hline
  \rowcolor{udcpink!25}
  \textbf{Título de columna} & \textbf{Outro título de columna} & \textbf{Outro título máis} \\\hline
  \textit{Título de fila}    & Contido da cela                  & Contido da cela \\
  \textit{Título de fila}    & Contido da cela                  & Contido da cela \\
  \textit{Título de fila}    & \multicolumn{2}{c}{Contido da cela múltiple} \\
  \textit{Título de fila}    & Contido da cela                  & Contido da cela \\
  \end{tabular}
  \caption{Pé de táboa descritivo (táboa con celas que ocupan varias columnas)}
  \label{tab:exemplocolumnas}
\end{table}

Tamén pode resultar necesario facer o propio mais en varias filas da mesma columna,
como ocorre na táboa~\ref{tab:exemplofilas} (páxina~\pageref{tab:exemplofilas}).
Para isto é preciso o paquete \texttt{multirow}, incluído entre os recomendados no
ficheiro raíz do proxecto (\verb+memoria_tfg.tex+).

O uso de celas multifila requerirá do xuste da coloración das filas, a fin de manter
a coherencia entre o contido e o continente. Así, no canto de usar un único comando
\verb+rowcolors+ para indicar a alternancia en toda a táboa, usaremos o comando
\verb+rowcolor+ antes dunha fila que queiramos colorear, e o comando \verb+cellcolor+
dentro dunha cela que queiramos colorear.

\begin{table}[hp!]
  \centering
  \begin{tabular}{c|c}
  \rowcolor{udcpink!25}
  \textbf{Título de columna} & \textbf{Outro título de columna} \\\hline
  \multirow{2}{*}{\textit{Título de fila}} & \cellcolor{udcgray!25} Contido da cela \\
                                           & Contido da cela \\
  \rowcolor{udcgray!25}
  \textit{Título de fila}                  & Contido da cela \\
  \multirow{3}{*}{\textit{Título de fila}} & Contido da cela \\
                                           & \cellcolor{udcgray!25} Contido da cela \\
                                           & Contido da cela \\
  \rowcolor{udcgray!25}
  \textit{Título de fila}                  & Contido da cela \\
  \end{tabular}
  \caption{Pé de táboa descritivo (táboa con celas que ocupan varias filas)}
  \label{tab:exemplofilas}
\end{table}

Por suposto, pódense combinar nunha mesma táboa os dous tipos de celas (as que se 
estenden máis dunha fila e máis dunha columna), como na táboa~\ref{tab:exemplofilasecolumnas}
(páxina~\pageref{tab:exemplofilasecolumnas}).

\begin{table}[hp!]
  \centering
  \begin{tabular}{c|c|c}
  \rowcolor{udcpink!25}
  \multicolumn{3}{c}{\textbf{Cela en varias columnas}} \\\hline
  \rowcolor{udcpink!25}
  \textbf{Título de columna}               & \textbf{Outro título de columna}             & \textbf{Outro título máis} \\\hline
  \multirow{2}{*}{\textit{Título de fila}} & \cellcolor{udcgray!25} Contido da cela       & \cellcolor{udcgray!25} Contido da cela \\
                                           & Contido da cela                              & Contido da cela \\
  \rowcolor{udcgray!25}
  \textit{Título de fila}                  & \multicolumn{2}{c}{Contido da cela múltiple} \\
  \multirow{3}{*}{\textit{Título de fila}} & Contido da cela                              & Contido da cela \\
                                           & \multicolumn{2}{c}{\cellcolor{udcgray!25} Contido da cela múltiple} \\
                                           & \multicolumn{2}{c}{Contido da cela múltiple} \\
  \rowcolor{udcgray!25}
  \textit{Título de fila}                  & Contido da cela                              & Contido da cela \\
  \end{tabular}
  \caption{Pé de táboa descritivo (táboa con celas que ocupan varias columnas)}
  \label{tab:exemplofilasecolumnas}
\end{table}

\section{Inclusión de código fonte}

Se precisamos incluír fragmentos de código fonte, podemos facelo, por exemplo, da
seguinte maneira:

\begin{lstlisting}[language=C]
#include <stdio.h>
#define N 10

int main()
{
  int i;

  // Isto é un comentario
  puts("Ola, mundo!");

  for (i = 0; i < N; i++)
  {
    puts("LaTeX é a ferramenta de edición ideal para profesionais da informática!");
  }

  return 0;
}
\end{lstlisting}

Para usar o estilo anterior en palabras ou expresións curtas que vaian insertadas nun
parágrafo de texto, recoméndase facelo \lstinline{deste_xeito(empregando, este, comando)}.

\section{Uso da relación de acrónimos e do glosario}

Os acrónimos edítanse no ficheiro \texttt{bibliografia/acronimos.tex}
e úsanse empregando a orde \texttt{acrlong} para obter o termo
completo (deste xeito: \acrlong{erlang}), a orde \texttt{acrshort}
para obter o acrónimo (deste xeito: \acrshort{erlang}). A primeira vez
que usamos un termo con acrónimo no documento é recomendable usar orde
\texttt{acrfull} (que produce ambas versións á vez:
\acrfull{erlang}). Os acrónimos que non se usan no documento, non
aparecen na relación que se xerar na versión PDF.

Pola súa banda, os termos do glosario edítanse no ficheiro
\texttt{bibliografia/glo\-sa\-rio.tex} e úsanse empregando a orde
\texttt{gls} (deste xeito, \gls{bytecode}) ou \texttt{Gls} (deste
xeito, \Gls{bytecode}). Ao igual que os acrónimos, os termos que non
se usan no documento, non aparecen na relación que se xera na versión
PDF.

%\include{contido/...}
%\include{contido/conclusions}
 \chapter{Introducción y motivación del dominio de búsqueda} 
\label{chap:introduccion} 
Este capítulo presenta la motivación y los objetivos del proyecto.
Viajar permite conocer nuevos lugares, culturas, tradiciones y formas de vida.
Sin embargo, la preparación de un viaje puede requerir la consulta de numerosas 
páginas web para encontrar información sobre destinos, rutas, actividades o puntos 
de interés. A partir de esta necesidad se plantea el desarrollo de un buscador 
especializado en viajes por España, diseñado para reunir información turística 
procedente de distintas fuentes y facilitar su consulta.

\section{Motivación}
La elección del dominio está estrechamente relacionada con la composición y los 
intereses del grupo. Sus integrantes proceden de diferentes nacionalidades, lo que
 aporta una visión multicultural y distintas experiencias en torno a los viajes. 
 Al mismo tiempo, todos comparten un gran interés por descubrir lugares, conocer
  nuevos destinos y acceder a información que ayude a organizar futuras visitas.

Esta afinidad llevó al equipo a seleccionar un ámbito que resultase motivador y cercano,
 pero que también tuviese utilidad práctica. El turismo ofrece una gran variedad de 
 contenidos sobre ciudades, regiones, rutas, naturaleza, gastronomía y patrimonio. No 
 obstante, estos contenidos suelen estar repartidos entre diferentes portales y guías,
  por lo que un sistema de búsqueda específico puede ayudar a localizarlos de forma más rápida y ordenada.

De este modo, el proyecto combina la motivación personal de un equipo multicultural
interesado en viajar con el propósito académico de aplicar técnicas de recopilación,
tratamiento, indexación y búsqueda de información. El resultado pretende ser una herramienta
útil para explorar destinos españoles y obtener información relevante de una manera centralizada y accesible.

Algunos ejemplos de consulta que podrían realizarse con el buscador incluyen:
\begin{itemize}
    \item Lugares que visitar en Santiago de Compostela.
    \item Rutas de naturaleza en Andalucía.
    \item Gastronomía y enoturismo en La Rioja.
    \item Playas de las Islas Baleares.
\end{itemize}

\section{Objetivos} 
El objetivo general consiste en desarrollar un buscador especializado en información
turística sobre España. Los objetivos específicos son los siguientes: 
\begin{enumerate} 
\item Obtener documentos de diferentes fuentes mediante crawling. 
\item Limpiar y normalizar los datos recopilados. 
\item Indexar los documentos mediante Elasticsearch. 
\item Diseñar consultas adaptadas al dominio turístico. 
\item Desarrollar una interfaz web para realizar búsquedas. 
\end{enumerate} 
\section{Estructura de la memoria} 
La memoria se organiza en cinco capítulos. Después de esta introducción, se presentan
el dominio y las fuentes de datos. A continuación, se describe el desarrollo del sistema,
incluyendo el crawling, la preparación de los datos y la indexación. Finalmente, se explican
las decisiones de búsqueda y las funcionalidades de la aplicación. 
 \chapter{Dominio y fuentes de Información} 
\label{chap:dominio_fuentes} 
\section{Alcance del Dominio}
El sistema se centra en la recuperación de contenido editorial en español relacionado con viajes por España. Se incluyen páginas dedicadas a destinos, ciudades, regiones, rutas y artículos temáticos sobre cultura, naturaleza, gastronomía, playas, deporte y otras propuestas turísticas. La unidad documental utilizada es una página web, de forma que cada documento del corpus representa una página de destino, una región, una ruta o un artículo turístico.

El objeto del buscador no es gestionar todos los aspectos relacionados con la contratación de un viaje. Por ese motivo, quedan fuera del alcance:

La consulta de precios de vuelos y alojamientos.
La disponibilidad en tiempo real.
La realización de reservas.
Las reseñas y puntuaciones publicadas por usuarios.
La comparación comercial entre hoteles, vuelos o paquetes turísticos.

Esta delimitación permite trabajar con información relativamente estable y de carácter descriptivo. Los precios y la disponibilidad cambian según las fechas, mientras que las reservas requieren integración con sistemas externos y operaciones transaccionales que no forman parte de los objetivos del proyecto. Las reseñas de usuarios tampoco se incluyen porque constituyen contenido generado por terceros y pueden incorporar datos personales o información de calidad difícil de controlar.

En consecuencia, el sistema no pretende sustituir a una plataforma comercial de reservas. Su objetivo es ayudar al usuario durante la fase previa de búsqueda y descubrimiento, proporcionando información que facilite la elección de destinos y actividades.

\section{Selección de Fuentes de Información}
Para construir la colección documental se seleccionaron dos fuentes:

Spain.info, portal oficial de turismo de España.
Wikivoyage en español, guía de viajes colaborativa.

La selección de dos sitios diferentes responde a los requisitos académicos del proyecto, que contemplan la obtención de documentos desde, al menos, dos sitios web y su preparación para la indexación en Elasticsearch.

Ambas fuentes ofrecen contenidos relacionados con el mismo dominio, pero presentan enfoques y estructuras diferentes. Esta complementariedad permite construir una colección más variada y obliga a desarrollar un proceso específico de rastreo y extracción para cada sitio.

\subsection{Spain.info}
Spain.info es el portal oficial de turismo de España y constituye la fuente institucional del proyecto. El sitio proporciona información sobre destinos, regiones, ciudades, rutas y artículos relacionados con patrimonio cultural, naturaleza, gastronomía, playas, deporte y otras experiencias turísticas.
La URL semilla utilizada para comenzar el rastreo es:
\url{https://www.spain.info/es/}
A partir de esta dirección se siguen enlaces pertenecientes a secciones relevantes para el dominio:

/es/destino/
/es/region/
/es/ciudades/
/es/arte-cultura/
/es/costas-playas/
/es/naturaleza/
/es/gastronomia-enoturismo/
/es/deporte-aventura/
/es/rutas/

Se excluyen las páginas de información práctica, los enlaces externos, los fragmentos de URL y las páginas que no aportan contenido temático al buscador.

La elección de Spain.info resulta adecuada porque ofrece información institucional, estructurada y centrada específicamente en el turismo nacional. Además, las comprobaciones iniciales realizadas con scrapy shell mostraron que la portada respondía con un código HTTP 200, proporcionaba alrededor de 196 KB de HTML y contenía 156 enlaces, entre los que se encontraban rutas a destinos, regiones y secciones temáticas. El contenido principal estaba disponible en el HTML descargado, por lo que no era necesario ejecutar JavaScript para extraerlo.

\subsection{Wikivoyage}
Wikivoyage es una guía de viajes colaborativa basada en MediaWiki. Sus páginas ofrecen información extensa sobre destinos y suelen organizarse en secciones como «Ver», «Hacer», «Comer» o «Dormir». Esta estructura proporciona una cantidad considerable de contenido textual y complementa el enfoque institucional de Spain.info.
La URL semilla seleccionada es:
\url{https://es.wikivoyage.org/wiki/Portal:España}
El rastreo sigue enlaces internos que comienzan por /wiki/. Se excluyen las anclas y las páginas especiales o pertenecientes a otros espacios de nombres, identificadas generalmente por la presencia del carácter : en la URL.
Wikivoyage se seleccionó porque sus contenidos se encuentran disponibles como HTML generado por MediaWiki y porque sus artículos ofrecen texto detallado sobre lugares, actividades y recomendaciones. Además, su carácter colaborativo introduce una estructura y un estilo diferentes a los de la fuente institucional, aumentando la diversidad de la colección.

El contenido de Wikivoyage se publica bajo licencia CC BY-SA, por lo que el sistema conserva la procedencia y la URL original de cada documento para mantener la correspondiente atribución.

\section{Complemento de las Fuentes}
La elección conjunta de Spain.info y Wikivoyage permite integrar dos perspectivas complementarias:
\begin{table}[htbp] 
    \centering 
    \caption{Comparación de las fuentes de información seleccionadas} 
    \label{tab:comparacion-fuentes} 
    \begin{tabular}{|p{3cm}|p{4cm}|p{5cm}|} 
        \hline \textbf{Aspecto} & \textbf{Spain.info} & \textbf{Wikivoyage} \\ 
        \hline Naturaleza de la fuente & Institucional & Colaborativa \\ 
        \hline Enfoque & Promoción turística oficial & Guías prácticas y descriptivas \\ 
        \hline Contenido principal & Destinos, regiones, rutas y artículos temáticos & Guías extensas de destinos \\ 
        \hline Estructura & Páginas organizadas por categorías turísticas & Artículos divididos en secciones \\ 
        \hline Información destacada & Descripción institucional y clasificación territorial & Información práctica y contenido textual amplio \\ 
        \hline 
    \end{tabular} 
\end{table}
La combinación de estas fuentes mejora la variedad del corpus. Spain.info aporta contenidos institucionales y una clasificación temática clara, mientras que Wikivoyage proporciona textos más extensos y secciones útiles para identificar actividades o aspectos concretos de un destino.

Las diferencias estructurales también representan un reto técnico relevante. Cada sitio necesita un spider específico, ya que las rutas, la estructura HTML y los selectores utilizados para extraer la información son diferentes. No obstante, los documentos obtenidos se transforman posteriormente a un esquema común, lo que permite almacenarlos e indexarlos de manera uniforme. 
 \include{contido/03_desarrollo} 
 \include{contido/04_decisiones_busqueda} 
 \include{contido/05_funcionalidades}

 %%%%%%%%%%%%%%%%%%%%%%%%%%%%%%%%%%%%%%%%
 % Apéndices, glosarios e bibliografía  %
 %%%%%%%%%%%%%%%%%%%%%%%%%%%%%%%%%%%%%%%%

%\appendix
%\appendixpage
%\include{anexos/adicional}
%\include{anexos/...}

% \printglossary[type=\acronymtype,title=\nomeglosarioacronimos]
%\printglossary[title=\nomeglosariotermos]

 \bibliographystyle{IEEEtranN}
 \bibliography{\bibconfig,bibliografia/bibliografia}
 \clearpage
 
\end{document}

%%%%%%%%%%%%%%%%%%%%%%%%%%%%%%%%%%%%%%%%%%%%%%%%%%%%%%%%%%%%%%%%%%%%%%%%%%%%%%%%
